\section{Evaluation Plan}

The first evidence layer is fixture correctness: path graphs, two-clique graphs,
determinism, and recursive split handling. The next evidence layer should
compare spectral construction against METIS, GeoSection, and clustering on
selected real-data smoke cases.

\begin{table}[h]
\centering
\begin{tabular}{llll}
\toprule
Claim & Evidence now & Missing evidence & Status \\
\midrule
Deterministic fixture output & L0 path/two-clique tests & None for fixtures & Current \\
Odd-\(k\) recursion handling & Recursive split fixtures & Broader \(k\) sweep & Current \\
RPLAN audit lineage & L1 sidecar checks & Real export transcript & Current \\
Quality vs. METIS/GeoSection & None in this paper & Smoke commands and outputs & Future \\
Runtime scaling & None in this paper & Benchmark script and hardware & Future \\
\bottomrule
\end{tabular}
\caption{Evidence ladder for T.14 claims. Fixture correctness, audit checks,
and empirical plan quality are separate evidence levels.}
\end{table}

\begin{table}[h]
\centering
\begin{tabular}{llll}
\toprule
Constructor & Primary role & Exposed failure mode & Audit surface \\
\midrule
Spectral & deterministic baseline & weak cut or imbalance & lineage and validation \\
METIS & multilevel baseline & dependency and heuristic variance & manifest lineage \\
GeoSection & geographic baseline & geometric oversimplification & assignment checks \\
Capacity clustering & capacity baseline & infeasible capacity/repair & status sidecars \\
Regionalization & merge baseline & merge-order sensitivity & merge witnesses \\
\bottomrule
\end{tabular}
\caption{Planned constructor comparison framing. Quantitative entries remain
future work until commands, data, and outputs are recorded.}
\end{table}
